\documentclass[12pt, a4paper]{article}

\usepackage[T1]{fontenc}
\usepackage[utf8]{inputenc}
\usepackage[english]{babel}
\usepackage[margin=2.5cm]{geometry}
\usepackage{xcolor}
\usepackage{hyperref}
\usepackage{parskip}
\usepackage{setspace}
\usepackage{microtype}
\usepackage{amsmath}
\usepackage{booktabs}
\usepackage{array}
\usepackage{longtable}

\definecolor{reviewerblue}{RGB}{0,0,180}
\definecolor{notegreen}{RGB}{0,120,0}

\hypersetup{colorlinks=true, linkcolor=reviewerblue, urlcolor=reviewerblue}

% -----------------------------------------------------------
\begin{document}
\setstretch{1.15}

% ============================================================
%  COVER LETTER
% ============================================================
\noindent
\textbf{Original Manuscript ID:} Access-2026-37050

\medskip
\noindent
\textbf{Original Article Title:} Artificial Intelligence Applications in International Large-Scale
Assessments: A Survey with LLM-Assisted Evidence Synthesis

\bigskip
\noindent
To: IEEE Access Editor

\medskip
\noindent
Re: Response to reviewers

\bigskip
\noindent
Dear Editor,

\medskip
Thank you for allowing a resubmission of our manuscript, with an opportunity to address the
reviewers' comments. We are uploading (a) our point-by-point response to the comments (below)
under ``Author's Response Files'', (b) an updated manuscript with yellow highlighting indicating
changes (as ``Highlighted PDF''), and (c) a clean updated manuscript without highlights (``Main
Manuscript'').

\medskip
\noindent
Best regards,\\[6pt]
Merve Dede and Ekrem \c{C}etinkaya

\bigskip
\noindent\textit{Note on location references.} The revised manuscript employs continuous line
numbering via the \LaTeX{} \texttt{lineno} package. All page and line references below correspond
to the visible line numbers in the compiled ``Main Manuscript'' PDF. Because these numbers differ
from LaTeX source-file line counts, readers should verify each citation directly against the
submitted PDF.

\newpage

% ============================================================
%  PRELIMINARY NOTE
% ============================================================

\noindent Dear Editor and Reviewers,

\medskip
Thank you for the careful and constructive evaluation of our manuscript. We have addressed every
comment in full and revised the manuscript accordingly. For each response, the corresponding
revision is identified by its section number and, where unambiguous, by the page and line numbers
visible in the compiled revised manuscript PDF. We note that the original reviewer letter contains
a numbering inconsistency (two comments labelled ``Comment \#3'' and no ``Comment \#2'' for
Reviewer~1); we follow the sequential ordering in which the comments appear.

\bigskip\hrule\bigskip

% ============================================================
%  REVIEWER 1
% ============================================================

\noindent\textbf{\large Reviewer \#1}

\bigskip

% -----------------------------------------------------------
\noindent\textbf{Comment \#1}

\medskip
\noindent\textcolor{reviewerblue}{%
The review relies primarily on Web of Science (WoS) as a main source of papers, supplemented by
Google Scholar searches and backward snowballing. This choice may result in somewhat incomplete
coverage. The authors should either broaden the database coverage or provide a substantially
stronger justification for the current approach. The exact Boolean search strings and field
restrictions should also be reported to ensure reproducibility.}

\medskip
\noindent\textbf{Response:}

Thank you for this constructive observation. We have substantially revised Section~3.2 (Literature
Review Strategy) to address both concerns.

\emph{Database justification.} We attempted an IEEExplore search using the identical query; it
returned no results due to that platform's field-length constraints. We note that the majority of
IEEE-indexed journals within this review's scope are indexed in WoS, so our primary search already
captures most of this literature; coverage of IEEE conference proceedings via WoS is acknowledged
as a minor limitation and is now stated explicitly. Google Scholar was retained as a complementary
source to reduce coverage bias, particularly for interdisciplinary conference proceedings, preprints,
and cross-disciplinary outlets at the intersection of AI and educational measurement.

\emph{Search-string transparency.} The full Boolean query is now reported in Section~3.2. The
analytical-methodology block was organised into four named subgroups from the outset: (B)~core
AI/ML methods, (C)~explainability methods (including ``explainable AI,'' ``XAI,'' ``SHAP''),
(D)~educational data methods (including ``educational data mining,'' ``learning analytics,''
``process data''), and (E)~psychometric-computational methods. With this expanded and fully
disclosed query, the WoS yield increased from 675 to 991 records; the final corpus grew from 130
to 212 studies. The complete pipeline scripts remain publicly available via the Data Availability
Statement.

\medskip
\noindent\textbf{Revision in the Manuscript:}
Section~3.2 was revised to include: the complete five-block Boolean query (Blocks A--E) displayed
verbatim; an explicit justification for WoS as the primary database; a description of the
IEEExplore attempt and its outcome; and a clarification of Google Scholar's role alongside an
acknowledgement of its limitations.

\medskip
\noindent\textbf{Location:} Section~3.2 (Literature Review Strategy). In the compiled PDF, the
Boolean query display begins on the page containing Section~3.2 and spans the passage running from
``The search query combined two conceptual blocks'' through ``Block A AND (Block B OR Block C OR
Block D OR Block E).''

\bigskip\hrule\bigskip

% -----------------------------------------------------------
\noindent\textbf{Comment \#2}

\medskip
\noindent\textcolor{reviewerblue}{%
The reported search terms do not clearly include several concepts that later form major parts of
the review, such as ``educational data mining'', ``learning analytics'', ``explainable AI / xAI'',
just to name a few. The authors should clarify whether these terms were included in the complete
search strategy and discuss the potential implications for study retrieval.}

\medskip
\noindent\textbf{Response:}

Thank you for this important observation. We confirm that these terms were included in the original
search strategy and were \emph{not} added post-hoc. However, the original manuscript did not
report the full Boolean string explicitly, which created the impression that these concepts were
absent. This has been corrected in the revised Section~3.2.

The revised query explicitly documents Block~C (explainability methods: ``explainable AI,'' ``XAI,''
``interpretable machine learning,'' ``SHAP,'' ``feature importance'') and Block~D (educational data
methods: ``educational data mining,'' ``learning analytics,'' ``process data,'' ``log-file
analysis,'' ``response process data'') as integral subgroups. The revised text states explicitly
that these subgroups were ``incorporated from the outset rather than added post-hoc.'' The expanded,
fully disclosed string yielded 991 WoS records and a final corpus of 212 studies.

\medskip
\noindent\textbf{Revision in the Manuscript:}
Section~3.2 was revised to display Blocks C and D verbatim, with an explicit statement that these
subgroups were incorporated from the outset.

\medskip
\noindent\textbf{Location:} Section~3.2. The relevant passage begins at ``Block C --- Explainability
methods'' and ``Block D --- Educational data methods'' in the query display.

\bigskip\hrule\bigskip

% -----------------------------------------------------------
\noindent\textbf{Comment \#3}

\medskip
\noindent\textcolor{reviewerblue}{%
The Policy Actionability Scoring Rubric is one of the manuscript's contributions, but the rubric
is currently internally defined and has not been externally validated. Please give more detail of
the scoring rubric.}

\medskip
\noindent\textbf{Response:}

Thank you for this important methodological point. Section~3.1.3 (Composite Actionability Scoring)
has been substantially revised in two respects.

\emph{Rubric detail.} The three dimensions are now described with concrete criteria:
$D_1$~(inferential warrant) ranges from purely descriptive or synthesis claims at the lowest level
through correlational association to causal ML designs (e.g., Bayesian Causal Forests) at the
highest. $D_2$~(effect specification) distinguishes findings that quantify effect size (coefficients,
percentage-point differences, AUC gains) from directional-only or unquantified findings.
$D_3$~(population boundedness) assesses whether conclusions are constrained to a named, identifiable
population. Score-assignment rules for all five levels are now stated deterministically: Score~5
requires causal warrant with quantified, population-bounded findings; Score~4 requires correlational
warrant with quantified effect sizes and country-specific scope; Score~3 requires at minimum
directional empirical evidence; Scores~1--2 cover purely methodological studies, synthesis papers,
or findings that are globally unspecified. Table~2 presents the full rubric. The revised text also
acknowledges that the rubric was developed internally and has not been externally validated; it
should be interpreted as a structured analytical lens rather than a psychometrically calibrated
instrument.

\emph{Validation.} A human rater independently scored a randomly selected subsample of 37 studies
by applying the same rubric without access to the programmatic scores. Exact agreement was
86.5\% (32/37); Cohen's $\kappa=0.76$ (unweighted) and $\kappa_w=0.83$ (quadratic-weighted),
indicating substantial agreement on the ordinal scale. The five disagreements involved boundary
cases concerning population boundedness, effect quantification, and methodological classification
and are noted in Section~6 (Limitations).

\medskip
\noindent\textbf{Revision in the Manuscript:}
Section~3.1.3 now includes: (i)~full operational definitions of all three scoring dimensions with
illustrative examples; (ii)~explicit score-assignment rules for all five levels; (iii)~an
acknowledgement that GRADE/SIGN frameworks were not directly applicable and why; (iv)~the full
corpus-level score distribution ($\bar{x}=3.70$, $SD=0.80$, $n=212$); and (v)~the inter-rater
validation results.

\medskip
\noindent\textbf{Location:} Section~3.1.3 (Composite Actionability Scoring) and Table~2 (Policy
Actionability Scoring Rubric). The passage begins at ``Established evidence-grading frameworks
such as GRADE \ldots{} were not directly applicable'' and ends at the inter-rater reliability
statistics.

\bigskip\hrule\bigskip

% -----------------------------------------------------------
\noindent\textbf{Comment \#4}

\medskip
\noindent\textcolor{reviewerblue}{%
Regarding the interpretation of xAI. SHAP, LIME, and similar techniques improve the
interpretability of model predictions, but they do not by themselves make findings policy-defensible
or causal. The manuscript would benefit from a clearer discussion of the distinction between
explaining a model's prediction and explaining the underlying educational phenomenon.}

\medskip
\noindent\textbf{Response:}

Thank you for this conceptually important point. We fully agree that this distinction was not drawn
with sufficient clarity in the original manuscript.

The distinction is now stated explicitly in three locations. In Section~1 (Introduction), a sentence
was added noting that ``XAI methods explain model predictions rather than underlying educational
mechanisms; policy-defensible inference therefore requires combining XAI outputs with substantive
domain knowledge and, where possible, causal designs.'' In Section~3.3.3 (Explainable Artificial
Intelligence), an explicit statement appears: these methods ``explain model predictions rather than
underlying educational mechanisms; feature importance should therefore not be interpreted as causal
evidence.'' In Section~5.1 (Recommendations for Educators and Practitioners), practitioners are
directly cautioned: ``XAI can clarify model-specific predictor contributions but should not be
interpreted causally; instructional decisions should therefore integrate model outputs with
contextual knowledge.''

In addition, Section~3.3.3 includes a new paragraph documenting XAI heterogeneity across student
groups and national contexts, grounding the model-versus-phenomenon distinction in concrete empirical
evidence from the reviewed studies.

\medskip
\noindent\textbf{Revision in the Manuscript:}
The distinction between model-prediction explanations and educational-phenomenon explanations is
now stated explicitly in Sections~1, 3.3.3, and~5.1.

\medskip
\noindent\textbf{Locations:}
\begin{itemize}
  \item Section~1 (Introduction): the passage beginning ``XAI methods explain model predictions
        rather than underlying educational mechanisms.''
  \item Section~3.3.3 (Explainable Artificial Intelligence): the passage beginning ``Importantly,
        these methods explain model predictions rather than underlying educational mechanisms.''
  \item Section~5.1 (Recommendations for Educators and Practitioners): Recommendation~3 beginning
        ``XAI can clarify model-specific predictor contributions.''
\end{itemize}

\bigskip\hrule\bigskip

% -----------------------------------------------------------
\noindent\textbf{Comment \#5}

\medskip
\noindent\textcolor{reviewerblue}{%
The extraction pipeline presented in the manuscript is an interesting contribution, but schema
validation alone does not demonstrate extraction accuracy. Since the authors report full human
verification and preserve an audit trail, it should be possible to provide quantitative information
such as the proportion of fields requiring correction, field-level agreement, or common error
categories. This would substantially strengthen the methodological contribution.}

\medskip
\noindent\textbf{Response:}

Thank you for this constructive suggestion. We agree that reporting schema-validation pass rates
alone is insufficient, and we have added quantitative extraction-fidelity data to Section~3.1.2
(Reliability Safeguards).

The revised section reports fidelity at two levels. At the field-coverage level: across 22 core
fields and 212 studies (4,664 field-instances), 962 instances (20.6\%) received sentinel labels
(\textit{Not Reported}, \textit{N/A}, or empty), falling to 7.7\% (293/3,816) when survey-weight
and open-access metadata were excluded; the most frequently flagged substantive fields were
\textit{ml\_techniques} (37.3\%), \textit{sample\_size} (25.0\%), \textit{countries\_formatted}
(24.1\%), \textit{confounders} (17.9\%), and \textit{effect\_size} (17.0\%), reflecting
source-level non-reporting rather than extraction errors. At the primary-method level: LLM--human
agreement was 93.2\% (124/133) after semantic normalization across studies with non-missing values
in both sources; the nine disagreements arose from multi-method studies where the LLM identified
a different primary method rather than making a factual error. The revised text also clarifies that
because human edits were not systematically logged, sentinel and null counts provide a completeness
proxy rather than a direct measure of correction frequency, and that field categories most
frequently requiring correction were outcome fields, conclusion templates, and method-family
assignments.

\medskip
\noindent\textbf{Revision in the Manuscript:}
Section~3.1.2 (Reliability Safeguards) now includes field-level sentinel rates, the primary
method-agreement statistic (93.2\%), a description of the nine disagreement cases, and an
identification of the field categories most frequently requiring human correction.

\medskip
\noindent\textbf{Location:} Section~3.1.2 (Reliability Safeguards). The relevant passage begins
at ``Extraction fidelity was assessed at field and method-agreement levels'' and continues through
the sentence ending ``a completeness proxy rather than a direct measure of correction frequency.''

\bigskip\hrule\bigskip

% -----------------------------------------------------------
\noindent\textbf{Comment \#6}

\medskip
\noindent\textcolor{reviewerblue}{%
The manuscript in current form is indeed very information-dense. The extensive descriptive material
on the structure of individual OECD and IEA assessments, particularly Tables~1--4, could be
condensed. This would allow the manuscript to focus more strongly on its main contributions.}

\medskip
\noindent\textbf{Response:}

Thank you for this suggestion. The four tables covering OECD and IEA assessment frameworks have
been replaced with a single consolidated summary table (Table~1) retaining only the dimensions
directly relevant to understanding how these assessments generate the data analyzed across the
reviewed studies: target group, core domains, cycle frequency, administration mode, psychometric
model, and scoring approach. Descriptive prose in Section~2.2 has been shortened accordingly.
This restructuring reduces the descriptive weight of Section~2 and rebalances the manuscript so
that the analytical contributions in Sections~3--5 receive proportionally greater emphasis.

\medskip
\noindent\textbf{Revision in the Manuscript:}
Sections~2.2 was revised by consolidating four detailed tables into Table~1 (``Overview of OECD
and IEA ILSAs: core characteristics and assessment framework'').

\medskip
\noindent\textbf{Location:} Section~2.2 (International Large-Scale Assessments) and Table~1.

\newpage

% ============================================================
%  REVIEWER 2
% ============================================================

\noindent\textbf{\large Reviewer \#2}

\bigskip

% -----------------------------------------------------------
\noindent\textbf{Comment \#1}

\medskip
\noindent\textcolor{reviewerblue}{%
Even though the paper is well written, the literature survey is weak and the introduction is
insufficient to capture the main objective of this work.}

\medskip
\noindent\textbf{Response:}

We thank the reviewer for this general concern. We note that this comment does not identify a
specific deficiency; our response therefore addresses the two dimensions most plausibly implied:
scope of the literature survey and clarity of the introduction's objectives.

\emph{Introduction.} The revised Introduction now states the study's primary objective, its scope
(212 studies, 2020--August 2026, seven ILSA programs), and the three specific departures from prior
syntheses more explicitly. The final paragraph of Section~1 identifies the five-category framework,
the LLM-assisted extraction contribution, and the policy actionability framework as the three
substantive contributions, and the sentence structure of that paragraph was revised to make the
motivation-to-contribution logic explicit.

\emph{Literature survey.} The survey is organized around the five-category framework in Section~4
and grounded in the methodological typology in Section~3.3. The revised manuscript covers 212
studies (up from 130), and the temporal scope was extended to August~2026. While we acknowledge
that a comprehensive discussion of every subfield covered would extend the manuscript beyond
reasonable length, we believe the revised Introduction and the explicit five-category framework in
Section~4 more clearly communicate the manuscript's scope and objectives.

\medskip
\noindent\textbf{Revision in the Manuscript:}
Section~1 (Introduction) was revised to state the primary objective and the three contributions
more explicitly, with the final paragraph rewritten to articulate the motivation-to-contribution
logic.

\medskip
\noindent\textbf{Location:} Section~1 (Introduction), final two paragraphs.

\bigskip\hrule\bigskip

% -----------------------------------------------------------
\noindent\textbf{Comment \#2}

\medskip
\noindent\textcolor{reviewerblue}{%
Please provide the full meaning of acronyms when they first appear in the paper (e.g., PISA,
TALIS).}

\medskip
\noindent\textbf{Response:}

Thank you for this observation. The revised Introduction (Section~1) provides the full name of each
assessment program at first mention: Programme for International Student Assessment (PISA), Teaching
and Learning International Survey (TALIS), Programme for the International Assessment of Adult
Competencies (PIAAC), Trends in International Mathematics and Science Study (TIMSS), Progress in
International Reading Literacy Study (PIRLS), International Civic and Citizenship Education Study
(ICCS), and International Computer and Information Literacy Study (ICILS). The same convention has
been applied to all other technical acronyms (e.g., AI, EDM, LA, ML, DL, XAI) introduced
throughout the manuscript.

\medskip
\noindent\textbf{Revision in the Manuscript:}
All acronyms are defined at first use in Section~1, in accordance with IEEE Access style
requirements.

\medskip
\noindent\textbf{Location:} Section~1 (Introduction). The OECD program acronyms are defined in the
second paragraph of Section~1; IEA program acronyms appear in the same paragraph.

\bigskip\hrule\bigskip

% -----------------------------------------------------------
\noindent\textbf{Comment \#3}

\medskip
\noindent\textcolor{reviewerblue}{%
How was the mean actionability score of 3.2 achieved?}

\medskip
\noindent\textbf{Response:}

Thank you for raising this point. The original manuscript did not describe the scoring mechanism
with sufficient transparency. Section~3.1.3 has been substantially revised to explain this in full.

Each study receives a score from 1 to 5 based on three observable dimensions: $D_1$ (inferential
warrant), $D_2$ (effect specification), and $D_3$ (population boundedness). Scores are assigned
deterministically through a hierarchical rule set applied to the $(D_1, D_2, D_3)$ triple, with
no coder discretion. The corpus mean reported in the original manuscript ($\bar{x}=3.20$) reflected
the arithmetic average of deterministic scores across 130 studies. In the revised version, the
corpus has been expanded to 212 studies, yielding an updated mean of $\bar{x}=3.70$ ($SD=0.80$).
The full score distribution (Table~2) and the scoring rule set are reported in Section~3.1.3 and
are also available in the publicly accessible pipeline scripts.

\medskip
\noindent\textbf{Revision in the Manuscript:}
Section~3.1.3 and Table~2 now provide a complete explanation of how composite scores are derived
from the three dimensions, including the deterministic rule set and corpus-level summary statistics.

\medskip
\noindent\textbf{Location:} Section~3.1.3 (Composite Actionability Scoring) and Table~2 (Policy
Actionability Scoring Rubric).

\bigskip\hrule\bigskip

% -----------------------------------------------------------
\noindent\textbf{Comment \#4}

\medskip
\noindent\textcolor{reviewerblue}{%
Elaborate on Response Time Effort (RTE).}

\medskip
\noindent\textbf{Response:}

Thank you for this suggestion. The revised Section~3.3.2 (Deep Learning and Process-Data Analytics)
now includes a formal definition of RTE. RTE relates a student's observed response time to an
item-level normative time threshold derived from the empirical distribution of response times across
examinees. Values substantially below the normative level, commonly $\text{RTE}<0.20$, are used to
identify potentially disengaged responding consistent with random guessing; values near or above the
threshold are consistent with sustained engagement. The revised text states this definition explicitly
and clarifies that RTE operationalizes test-taking effort at the item level, complementing
accuracy-based indicators by capturing behavioral validity rather than correctness alone.

\medskip
\noindent\textbf{Revision in the Manuscript:}
Section~3.3.2 now includes a formal definition of RTE with its standard interpretation and threshold
criterion.

\medskip
\noindent\textbf{Location:} Section~3.3.2 (Deep Learning and Process-Data Analytics). The passage
begins at ``Response time effort (RTE), for instance, relates observed response time to an
item-level normative time threshold.''

\bigskip\hrule\bigskip

% -----------------------------------------------------------
\noindent\textbf{Comment \#5}

\medskip
\noindent\textcolor{reviewerblue}{%
Elaborate on how XAI heterogeneity manifests across student groups.}

\medskip
\noindent\textbf{Response:}

Thank you for this suggestion. The revised Section~3.3.3 (Explainable Artificial Intelligence) now
includes a dedicated passage on XAI heterogeneity at two levels.

At the student-group level, SHAP dependency plots from reviewed studies reveal that feature
contributions differ in direction and magnitude across achievement quartiles: for example,
self-efficacy shows diminishing marginal returns in the upper achievement quartile while carrying
the largest positive SHAP contribution in the lower quartile. At the cross-national level, the
relative importance ranking of the same variable shifts across systems even when the variable
remains statistically significant; ESCS dominates in resource-constrained contexts while
psychosocial constructs lead in higher-income systems. These patterns confirm that XAI findings
are context-bound inferences about model behavior rather than universal claims about educational
causation.

\medskip
\noindent\textbf{Revision in the Manuscript:}
Section~3.3.3 now includes an explicit passage on XAI heterogeneity across student achievement
groups and national contexts, supported by examples from reviewed studies.

\medskip
\noindent\textbf{Location:} Section~3.3.3 (Explainable Artificial Intelligence). The passage
begins at ``The reviewed studies further indicate heterogeneity in feature contributions across
student groups and national contexts.''

\bigskip\hrule\bigskip

% -----------------------------------------------------------
\noindent\textbf{Comment \#6}

\medskip
\noindent\textcolor{reviewerblue}{%
It is not clear how you obtained your mean composite actionability scores.}

\medskip
\noindent\textbf{Response:}

This comment addresses the same concern as Comment~\#3 above; the response and revision are
identical. For completeness: each study is scored deterministically from the $(D_1, D_2, D_3)$
triple using a hierarchical rule set (no coder discretion); the composite score is the resulting
ordinal value from 1 to 5; the corpus mean is the arithmetic average of these scores across all
212~studies ($\bar{x}=3.70$, $SD=0.80$). The scoring function is included in the publicly available
pipeline scripts. Full details are in the revised Section~3.1.3 and Table~2.

\medskip
\noindent\textbf{Location:} Section~3.1.3 and Table~2 (see Comment~\#3 above).

\bigskip\hrule\bigskip

% -----------------------------------------------------------
\noindent\textbf{Comment \#7}

\medskip
\noindent\textcolor{reviewerblue}{%
Provide separate sections for limitations and future directions.}

\medskip
\noindent\textbf{Response:}

Thank you for this structural suggestion. The combined ``Limitations and Future Directions''
subsection has been separated into two independent sections in the revised manuscript:
Section~6 (Limitations) and Section~7 (Future Directions). Each section now stands independently,
allowing readers to engage with methodological limitations and research priorities separately.

\medskip
\noindent\textbf{Revision in the Manuscript:}
Sections~6 and~7 are now distinct, labeled ``Limitations'' and ``Future Directions'' respectively.

\medskip
\noindent\textbf{Locations:}
\begin{itemize}
  \item Section~6 (Limitations): begins with ``Several methodological limitations warrant
        acknowledgement.''
  \item Section~7 (Future Directions): begins with ``Three research priorities emerge from the
        corpus.''
\end{itemize}

\bigskip\hrule\bigskip

% -----------------------------------------------------------
\noindent\textbf{Comment \#8}

\medskip
\noindent\textcolor{reviewerblue}{%
What are your recommendations for educators based on your findings in this work?}

\medskip
\noindent\textbf{Response:}

Thank you for this practically important suggestion. A dedicated subsection, Section~5.1
(Recommendations for Educators and Practitioners), has been added to the Conclusion.

Four concrete evidence-based recommendations are provided. (1)~Self-efficacy, metacognitive
strategy use, and school belonging are the most consistently modifiable predictors across programs
and geographies; school-level interventions targeting these constructs are more likely to generalize
than those targeting demographic or resource variables alone. (2)~Process data analytics reveal
that students reaching identical scores often follow measurably different behavioral pathways;
educators can use behavioral engagement indicators such as RTE and action-sequence profiles to
identify students who appear on-track by score alone but show signs of disengagement.
(3)~SHAP-based feature importance outputs should be interpreted as model-specific summaries of
predictor contributions rather than causal explanations, and should be triangulated with contextual
knowledge before informing instructional decisions. (4)~Given the cross-cultural instability of
predictors documented across the corpus, findings from high-performing or resource-rich systems
should be piloted rather than directly adopted in different national contexts.

\medskip
\noindent\textbf{Revision in the Manuscript:}
Section~5.1 (Recommendations for Educators and Practitioners) added to the Conclusion with four
evidence-based recommendations derived from the survey findings.

\medskip
\noindent\textbf{Location:} Section~5.1 (Recommendations for Educators and Practitioners). The
subsection begins at ``The findings yield four practical recommendations.''

\bigskip\hrule\bigskip

\newpage

% ============================================================
%  VERIFICATION TABLE
% ============================================================

\noindent\textbf{\large Summary Verification Table}

\medskip
\noindent\textit{Note:} Exact line numbers are not reproducible from the LaTeX source alone and
must be read directly from the compiled revised PDF. Section references below are stable across
compilation.

\medskip
\renewcommand{\arraystretch}{1.4}
\begin{longtable}{p{1.2cm} p{3.3cm} p{4.0cm} p{5.5cm}}
\toprule
\textbf{Rev.} & \textbf{Comment} & \textbf{Revision Made} & \textbf{Location} \\
\midrule
\endhead

R1\,\#1 & Database coverage and Boolean search string transparency
         & Full Boolean query (Blocks A--E) disclosed; IEEExplore attempt described; GS role
           clarified; corpus expanded to 212 studies
         & Section~3.2 (Literature Review Strategy) \\

R1\,\#2 & Missing EDM, LA, XAI search terms
         & Blocks C and D documented verbatim with explicit statement that subgroups were
           incorporated from the outset
         & Section~3.2 (Blocks C and D in query display) \\

R1\,\#3 & Scoring rubric detail and validation
         & Full operational definitions of $D_1$--$D_3$; five-level score rules; inter-rater
           agreement ($\kappa=0.76$, $\kappa_w=0.83$, $n_{\mathrm{sub}}=37$)
         & Section~3.1.3 and Table~2 \\

R1\,\#4 & XAI prediction vs.\ educational phenomenon distinction
         & Explicit clarification added in Sections~1, 3.3.3, and~5.1; new XAI heterogeneity
           paragraph in Section~3.3.3
         & Sections~1, 3.3.3, and~5.1 \\

R1\,\#5 & Quantitative extraction fidelity data
         & Sentinel rates (20.6\%), primary-method agreement (93.2\%), field-level correction
           categories reported
         & Section~3.1.2 (Reliability Safeguards) \\

R1\,\#6 & Condensing Tables 1--4 on OECD/IEA assessment structure
         & Four tables replaced with one consolidated Table~1
         & Section~2.2 and Table~1 \\

R2\,\#1 & Literature survey weak; introduction insufficient
         & Introduction revised to state objectives, scope, and contributions more explicitly
         & Section~1, final two paragraphs \\

R2\,\#2 & Acronyms not defined at first use
         & All program and technical acronyms defined at first mention in Section~1
         & Section~1 (Introduction) \\

R2\,\#3 & How was mean actionability score of 3.2 derived?
         & Deterministic rule set and corpus statistics ($\bar{x}=3.70$, $SD=0.80$, $n=212$)
           fully explained; rubric in Table~2
         & Section~3.1.3 and Table~2 \\

R2\,\#4 & Elaborate on RTE
         & Formal definition and threshold interpretation added
         & Section~3.3.2 \\

R2\,\#5 & XAI heterogeneity across student groups
         & Dedicated passage on student-group and cross-national heterogeneity added
         & Section~3.3.3 \\

R2\,\#6 & Composite actionability score derivation unclear
         & Same revision as R2\,\#3
         & Section~3.1.3 and Table~2 \\

R2\,\#7 & Separate sections for limitations and future directions
         & Sections~6 (Limitations) and~7 (Future Directions) created as independent sections
         & Sections~6 and~7 \\

R2\,\#8 & Recommendations for educators
         & Section~5.1 (Recommendations for Educators and Practitioners) added with four
           evidence-based recommendations
         & Section~5.1 \\

\bottomrule
\end{longtable}

\bigskip\noindent\textit{Note: References suggested by reviewers should only be added if relevant
to the article and make it more complete. Excessive cases of recommending non-relevant articles
should be reported to \href{mailto:ieeeaccesseic@ieee.org}{ieeeaccesseic@ieee.org}.}

\end{document}
